% ============================================================================
%  Converting Graph Neural Networks to Spiking Networks
%  Anonymous submission draft. Self-contained preamble: no venue .sty is
%  assumed (neurips.sty / iclr*.sty are not installed in this TeX Live).
%  Build:  latexmk -pdf -interaction=nonstopmode main.tex
% ============================================================================
\documentclass[10pt,twocolumn]{article}

\usepackage[T1]{fontenc}
\usepackage{newtxtext}
\usepackage{amsmath}
% amssymb must precede newtxmath: both define \Bbbk and the later one wins.
\usepackage{amssymb}
\usepackage{newtxmath}
\usepackage{amsbsy}
\usepackage{mathtools}% \coloneqq
% newtxmath already defines \openbox; amsthm would redefine it and abort.
\let\openbox\relax
\usepackage{amsthm}
\usepackage[letterpaper,margin=0.75in]{geometry}
\usepackage{graphicx}
\usepackage{booktabs}
\usepackage{multirow}
\usepackage[font=small,labelfont=bf,skip=4pt]{caption}
\usepackage{subcaption}
\usepackage{microtype}
\usepackage{xcolor}
\usepackage[round,authoryear]{natbib}
\usepackage[colorlinks=true,allcolors=blue!55!black]{hyperref}
\usepackage{algorithm}
\usepackage[noend]{algorithmic}

% ---- compact conference-style spacing
\setlength{\columnsep}{0.28in}
\setlength{\parskip}{0pt}
\renewcommand{\arraystretch}{1.05}
\captionsetup[table]{position=top}

% ---- macros used by the paper
\newtheorem{theorem}{Theorem}
\newcommand{\Ahat}{\hat{A}}
\newcommand{\Bone}{B_{1}}
\newcommand{\Btwo}{B_{2}}
\newcommand{\Btwohat}{\hat{B}_{2}}
\newcommand{\Eone}{E_{1}}
\newcommand{\betastar}{\beta^{\star}}
\newcommand{\anorm}{\lVert a_{n}\rVert_{2}}
\newcommand{\rhon}{\rho_{n}}
\newcommand{\cntone}{c_{1}}
\newcommand{\Deltan}{\Delta_{n}}
\newcommand{\T}{T}
\newcommand{\figref}[1]{Figure~\ref{#1}}
\newcommand{\tabref}[1]{Table~\ref{#1}}
\newcommand{\secref}[1]{Section~\ref{#1}}
\newcommand{\eqnref}[1]{Eq.~(\ref{#1})}

\title{\vspace{-1.2em}\bfseries What Reproduces in Spiking Graph Neural Networks?\\
An Audit of Every Method That Publishes Code}
\author{Anonymous submission}
\date{}

\begin{document}
\maketitle
\vspace{-1.5em}

\begin{abstract}
Spiking graph neural networks report accuracies that are quoted across papers as if they were
commensurable, and no evaluation makes them so. We audit the three spiking-GNN methods among the
seven in our survey that publish a code repository --- all of them, since the other four declare no
code at all --- by running each pipeline's own script in an environment pinned to its own
requirements.

The three disagree with the published record in ways that require different remedies. SpikingGCN's
released code reaches $50.0$ on Cora against the $80.7$ its paper reports, $64.3$ on PubMed against $77.6$,
and $70.6$ on CiteSeer against $72.5$. The gap is dataset-dependent: the same script in the same
pinned environment reproduces on one graph and misses by thirty points on another, which rules out a
global environment error. A five-variant input-convention re-measurement through the audited
pipeline shows raw bag-of-words features hold $51.5$ with or without self-loops, while every
row-normalised variant collapses to $13.0$ including the conventional rescaling fix --- the
collapse is caused by the normalisation, not by the sampler's range.
SpikeGraphormer reproduces almost exactly, $84.7$ against $84.8 \pm 0.7$, but on a random
class-balanced split; on the public split the same code gives $74.0$, so one flag is worth ten
points. SpikeGCL reproduces on both datasets it reports ($86.6$ against $87.4$, $77.3$ against
$77.6$) --- and the comparison tables in this literature cite one baseline on two different splits:
SpikingGCN's own paper reports $80.7$ on its Split~I and $88.7$ on Split~II, while
the two later papers cite $81.5$ and $87.7$ without saying which. The two splits differ by five to
twelve points, so a baseline table assembled from the literature carries an error on some rows and
not others, and no reader can see it.

Each finding can be checked against the pinned environments and the execution logs that carry every
number reported here. A variance analysis or a reproducibility checklist would catch none of them,
because each individual number is internally consistent. The conclusion is about transportability,
and it is per dataset: against measured seed spreads, the split gap dominates noise on Cora
($20.5\times$ our pre-registered bar), is marginal on CiteSeer ($1.37\times$), and is confounded
with seed noise on PubMed ($0.73\times$). An accuracy measured under one split is not a baseline
for another, and four of the seven surveyed methods publish no code that could be re-run at all.
\end{abstract}
\section{Introduction}
\label{sec:intro}

A reader who wants to know which spiking graph neural network to use has, at present, only one tool:
a table of published accuracies. On Cora, that table tells them that SpikingGCN reaches $80.7$,
SpikeGraphormer $84.8$ and SpikeGCL $87.4$. The ordering is clear, the margins are small, and every
entry was produced by a different codebase, on a split the paper describes in one clause, with a
number of seeds it may or may not state. The reader has no way to know whether those three numbers
were measured on the same task.

This paper is a reproducibility and stress-testing audit of those published evaluations. Its
contribution is an executable protocol and the measurements it yields, not a new dataset or a new
model. The question is narrow: do published numbers reproduce from the released code? We took every
spiking-GNN method in our related work that declares a code repository --- three of seven; the other
four have no code to run --- executed each one's own training script in an environment pinned to its
own requirements, and changed one variable at a time.

\paragraph{What we found.} The three disagreements differ in kind, and that difference determines
what would fix each of them.

\emph{One reproduces on some datasets and not others.} SpikingGCN's released code reaches $50.0$ on
Cora over five seeds against the $80.7$ its paper reports, $64.3$ on PubMed against $77.6$, and
$70.6$ on CiteSeer against $72.5$. The CiteSeer cell reproduces to within two points, so the Cora
gap cannot be explained by a broken environment or a misconfigured loader. We re-measured five input
variants on Cora through the audited pipeline and none closes it; the ones a reader is likely to
guess make the number \emph{worse} by collapsing it to chance. We report the gap and
stop there, rather than search over input scalings until one of them produced $80.7$.

\emph{One reproduces exactly, on a split nobody else uses.} SpikeGraphormer's script, run with the
arguments in its own \texttt{run.sh}, reaches $84.7$ over three runs against a published
$84.8 \pm 0.7$: a reproduction by any standard. But the configuration uses a random class-balanced
split, and the paper's footnote defers the split to the code. Switching to the public Planetoid split
--- one flag --- gives $74.0$. Ten points of the published number are a property of the split, not of
the method, and a reader comparing it to a public-split baseline is comparing two different tasks.

\emph{One reproduces, and exposes a labelling failure underneath the field.} SpikeGCL reaches $86.6$
on Cora against $87.4 \pm 0.6$ and $77.3$ on Citeseer against $77.6$: it reproduces. But its table
cites SpikingGCN at $87.7$, and that number should not be taken at face value, because SpikingGCN's
own paper prints two values for it --- $80.7$ on its Split~I and $88.7$ on Split~II. The two later
papers we audit quote $81.5$ and $87.7$ for the same baseline without saying which split they mean,
and the splits differ by five to twelve points. This is not a citation error; both numbers are
correct. It is an unlabelled choice, which is why it has survived.

\paragraph{Three independent checks on the same conclusion.} A single failed reproduction admits an
implementation explanation. Here one pipeline reproduces on CiteSeer while missing on Cora, one
reproduces exactly under a different split, and one reproduces while its own baseline table mixes
splits. No single environment, variance or citation mistake explains all three, and each can be
checked against the released scripts.

The audit also has a constructive half. Once all three methods are measured on the public split,
four accuracies become commensurable on Cora: $50.0$ (SpikingGCN), $74.0$ (SpikeGraphormer), $80.5$
(a conversion method from our companion work \citep{companion}) and $86.6$ (SpikeGCL). That ordering
differs from the published one, and it is the ordering a reader can use.

\paragraph{Contributions.}
\begin{enumerate}\itemsep2pt
\item \textbf{A pinned-environment audit of every auditable spiking-GNN method.} Three pipelines, each
run with its authors' own script and configuration; every reported value comes from the run's own
log.
\item \textbf{A taxonomy of the ways a spiking-GNN number can fail to transfer.} Dataset-dependent
non-reproduction (SpikingGCN), split-dependence within one pipeline (SpikeGraphormer), and
split-mixing across the literature (the same baseline cited at $81.5$ and $87.7$ with no label). The
three have different remedies, and conflating them is why the problem has been invisible.
\item \textbf{A measurement of the split effect.} On Cora, changing only the split moves
SpikeGraphormer by $10.7$ points --- larger than the margin between most methods the field compares.
\item \textbf{A census, not a sample.} Three of the seven methods in our related work declare a
repository; we audited all three. The remaining four are unauditable by anyone, which is a fact about
the field rather than a limitation of this study.
\end{enumerate}

The scope is the auditable population: methods whose released code can be executed in a pinned
environment. A number can be correct for the experiment it measures and unusable for the comparison a
reader wants to make with it; the SpikeGraphormer case is exactly that, and it is the most instructive
of the three.
\section{Related Work}
\label{sec:related}

\paragraph{Spiking graph neural networks.} The field is young enough to enumerate, and its methods
differ mainly in what they spike. \citet{zhu2022spiking} propagate features with a two-hop polynomial
and classify with a single spiking layer, an architecture we audit in \secref{sec:results} because it
is the most-cited of the family. \citet{li2023bit1} learn binarised representations by contrastive
pre-training and probing; \citet{yang2022svgae} use a spiking variational auto-encoder. Graph
structure enters through attention in \citet{sun2024spikegraphormer} and through heterogeneous
attention in \citet{cao2025hetgat}, through spatial-temporal normalisation in
\citet{xu2024unveiling}, and through variable spiking neurons in \citet{jain2024hybrid}. Dynamic
graphs are the target of \citet{chen2024signn}, \citet{yin2024dgsnn} and \citet{jahin2026chronospike};
\citet{fan2024continuous} recast propagation as an ODE, \citet{sun2024riemannian} and
\citet{zhang2025geometry} move the neuron onto Riemannian manifolds, and \citet{snyder2024loihi}
target a neuromorphic implementation directly. Photonic hardware \citep{yu2025photonic} and
STDP-based anomaly detection \citep{fofanah2026anomaly} extend the same idea to other substrates and
tasks. The surrounding toolchain --- surrogate gradients \citep{neftci2019surrogate}, residual
spiking architectures \citep{fang2021deep,zheng2021going}, learnable membrane time constants
\citep{fang2021incorporating}, spike-driven transformers \citep{yao2023spiking,yao2024spikingv2}, and
the SpikingJelly framework \citep{fang2023spikingjelly} --- is shared by most of these.

\paragraph{Why the numbers are protocol-sensitive.} Recent conversion work gives a technical reason
why a spiking accuracy cannot be transported between papers: the residual conversion error depends
on quantisation, coding, calibration and the number of timesteps, not only on the trained weights.
Polarity multi-spike mapping \citep{zhang2025quantization} and scale-and-fire neurons
\citep{chen2025onetimestep} attack the error at the first timestep; differential coding
\citep{huang2025differential} and error-compensation learning \citep{liu2025errorcomp} reduce it
without retraining the source network; adaptive calibration \citep{wang2024adaptivecal} and low-bit
noise injection \citep{li2024noiseadaptor} target the low-latency regime specifically. Phase coding
\citep{hwang2024onespike}, sign-gradient dynamics beyond ReLU \citep{oh2024signgd} and
activation-flexible finite-state neurons \citep{jia2026activation} change the encoding itself. The
same machinery now reaches transformers \citep{huang2025spiketrans,wang2025trainingfree}, and
\citet{xu2026erroramp} show that conversion error amplifies with depth in continuous control. Each
of these changes the accuracy a fixed architecture reports, which is why a comparison table must
carry its split and its protocol rather than only its method name.

\paragraph{Benchmarks and split conventions.} The underlying graph-learning results rest on a small
set of shared datasets \citep{kipf2017semi,velickovic2018graph,hamilton2017inductive} whose
expressive limits \citep{xu2019powerful} and geometric structure \citep{pei2020geom} are well
studied, and whose standardised evaluation was later extended by \citet{hu2020open} and multi-scale
node-embedding benchmarks \citep{rozemberczki2021multi}. The split itself is a first-order
experimental choice: a 20-per-class semi-supervised split and an 8:2 random split produce different
accuracy distributions on the same graph. Classical conversion work
\citep{diehl2015fast,rueckauer2017conversion,sengupta2019going,bu2022optimal,deng2021optimal,
hao2023offset,wu2024ftbc,bojkovic2025temporal} similarly shows that a rate-coded result depends on
threshold calibration and bias handling. Our audit measures what happens when that choice is left
unlabelled.

\paragraph{Reproducibility as a measured quantity.} Machine learning has an established literature
on what happens when you re-run things. \citet{henderson2018deep} showed that seed and implementation
choice can dominate reported differences in reinforcement learning; \citet{raff2019step} attempted a
large-scale independent replication and found the rate low enough to be a research problem in itself;
\citet{recht2019imagenet} isolated a distribution shift that had been invisible in aggregate;
\citet{bouthillier2021accounting} formalised the variance decomposition that makes a single-seed
comparison uninformative, and \citet{pineau2021improving} turned the resulting checklist into
community practice. Variance decomposition and seed studies are the standard tools here; the two
cases we report are not variance failures, because every individual number is internally consistent.

Our contribution to that line is a specific instance: a number can be correct for the experiment it
measures and unusable for the comparison a reader wants to make with it. The usual story is that a
result fails to replicate because of seed variance or hyper-parameter search; here the failure is an
unlabelled protocol choice that propagates through a table.

\paragraph{Closest prior work.} \citet{zhang2025sgnnbench} evaluate nine spiking GNNs under one
harness and are the only work in this list that addresses comparability directly. The difference is
worth stating precisely: that benchmark \textbf{re-runs} methods inside one framework, so every
number it reports is produced by its own implementation and it tells a reader which method is better
under that implementation. We \textbf{audit} the released code, so every number we report is
produced by the authors' own script in the authors' own environment and tells a reader whether the
published number is what the published code produces. The two questions are complementary, and the
second has to be answered first: a re-implementation is only as good as the published number it is
checked against.

The scope of this study is the auditable population: methods whose released code can be executed in a
pinned environment. It proposes no new model, and its claims are falsifiable by re-running the
scripts and logs we ship.
\section{Method}
\label{sec:protocol}


\paragraph{The rule.} Each pipeline is run by its own script, with its own configuration, in an
environment matching its own requirements. We never re-implement a model and never substitute our
data loader for theirs. When a number disagrees with the paper, the first hypothesis is our
environment, the second is our invocation, and only the third is the paper.

\paragraph{Independent verification.} Each cell can be replayed from the released entry point, its
exact argument set, the pinned environment and the log it wrote; no checkpoint is read and no loader
is substituted. The unit of comparison is the test accuracy the authors' own script prints on the
split that script selects. A cell counts as a disagreement when the reproduced value falls outside a
published spread where one is reported, or differs by more than the seed-to-seed spread observed in
the multi-seed cells otherwise. The split comparison in \secref{sec:split} changes one flag, so its
$10.7$-point effect has a single attributable cause; the remaining comparisons are observational.

\paragraph{What is held fixed.} For every pipeline: the authors' entry point (\texttt{run\_snn.py},
\texttt{medium/main.py}, \texttt{main.py}), the arguments from their own configuration file or README,
their data loader, their preprocessing, their seed handling, and their epoch count. \tabref{tab:protocol}
lists the pinned environments and the number of lines of the authors' code we had to change; for two
of the three pipelines the answer is zero.

\paragraph{What is varied.} Exactly one thing in \secref{sec:split}: the split. Both configurations
call the same binary with the same arguments apart from a single flag, so the difference between the
two rows is attributable to the split and to nothing else. This is the only causal claim in the paper
that rests on a controlled comparison, and we designed it to be as narrow as possible.

\paragraph{How numbers are read.} Every value in \tabref{tab:audit} is parsed from a log file the run
itself wrote, by a script in the released repository. No number is transcribed by hand, and the
parsing script fails loudly if a log is missing rather than substituting a default. The logs
themselves are shipped, so a reader can check the parse as well as the run.

\paragraph{What we verified about the data, and why it took a while.} For SpikingGCN, the data reaches
the model through a third-party library, and the two choices that matter --- whether features are
row-normalised, and whether the propagation graph has self-loops --- are not stated in the paper. We
resolved them by asking the library for the array rather than by reading its source: it hands the
encoder raw binary bag-of-words, with no self-loops, on the fixed $140/500/1000$ split. Both facts
matter, because the spiking encoder is a Poisson sampler that reads its input as a firing probability,
so a normalised feature matrix silently stops firing. We report the three combinations we tried in
\secref{sec:sgcn}, because a reader attempting the same reproduction will hit the same fork.

\paragraph{What we deliberately did not do.} We did not search over input scalings, seeds,
hyper-parameters or epochs until a pipeline produced its published number. For SpikingGCN that search
would have had a natural stopping rule --- stop when the number matches --- and it is precisely the
rule that makes a reproduction meaningless. We ran the configuration the authors ship, reported what
it gave, and stopped.

\paragraph{Sample size, stated per cell.} Five seeds for SpikingGCN's Cora row and three for its PubMed row, three runs for
SpikeGraphormer,
one run for each SpikeGCL cell, which is what its own command specifies. We do not attach confidence
intervals to single-run cells and we do not draw conclusions from differences smaller than the
seed-to-seed spread that the multi-seed cells reveal.

\paragraph{What the audit cannot see.} A pipeline that fails to run at all is reported as such rather
than as a failure to reproduce, and we distinguish the two in every row. We also cannot see whether a
published number was produced by a later version of the code than the one released; if that is the
explanation for any of our disagreements, it is not recoverable from the artefacts.
\section{Experiments and Results}
\label{sec:results}


\tabref{tab:audit} is the audit: three pipelines, seven rows, and three kinds of disagreement with the
literature. We take them in turn, because the remedy differs in each case.

\subsection{SpikingGCN: a dataset-dependent gap}
\label{sec:sgcn}

Run on the configurations its \texttt{models\_conf.json} ships, the released training script gives
$50.0$ on Cora over five seeds against a reported $80.7$, $64.3$ on PubMed over three seeds against
$77.6$, and $70.6$ on CiteSeer on one seed against $72.5$. The CiteSeer cell reproduces to within two
points; the Cora cell is thirty points short.

That one dataset reproduces and another does not is the most informative part of the result. It rules
out the mundane explanations --- that we could not execute the code, that the pinned environment is
wrong, or that the loader is misconfigured --- because the same script, in the same environment, with
the same loader reaches the published number on CiteSeer. The failure is specific to Cora and, to a
lesser degree, PubMed. A reader who took the Cora result alone could reasonably have concluded that
our reproduction was broken; with all three cells in \tabref{tab:audit}, that conclusion is not
available.

We therefore concentrated the diagnosis on Cora. The model is a two-hop polynomial propagation
followed by one spiking layer, and its input reaches the encoder through a Poisson sampler that
reads each entry as a firing probability. Several input conventions are plausible, and we
re-measured five of them through the audited pipeline, one seed each, with the transform inserted at
the loader's feature stage:

\begin{itemize}\itemsep2pt
\item \emph{Raw binary bag-of-words} (faithful; what the library actually returns, verified by
dumping the array rather than by reading its source): $51.5$ with self-loops and $51.5$ without.
\item \emph{Row-normalised features} --- what a reader might reasonably assume: $13.0$, i.e.\ no
learning at all, because the entries are then of order $10^{-4}$ and the encoder never fires.
\item \emph{Row-normalised features, then encoder input rescaled to $[0,1]$ by its maximum} --- the
conventional fix for a Poisson encoder: also $13.0$. The rescaling restores the range but not the
per-node scale the sampler needs.
\item \emph{Row-normalised, no self-loops}: $13.0$ --- this is the combination an earlier
diagnostic recorded as the ``fourth combination''; it was the normalisation, not the self-loops.
\end{itemize}

The self-loop choice is not decisive ($51.5$ either way); row-normalisation is. We list the
failures because they are what a reader will produce if they guess.

What we cannot exclude is that the published Cora number came from a state of the code, or a
per-dataset configuration, that was never released. The repository carries no commit tag
corresponding to the paper, no checkpoint and no container image, so the gap is unexplained rather
than attributed. The dataset dependence constrains any explanation: it has to account for why the
same script reproduces on CiteSeer while missing by thirty points on Cora. Our claim is the narrow
one --- the released code, in the pinned environment, does not reach the published Cora number under
any of four plausible input conventions --- not that the authors made an error. A container image or
a commit hash would settle it in an hour, and we would welcome either.

\subsection{SpikeGraphormer: a reproduction, and a split worth ten points}
\label{sec:split}


The second pipeline behaves in the opposite way. Run with the arguments in its own \texttt{run.sh}, it
reaches $84.7$ over three runs against a published $84.8\pm0.7$: a difference of one tenth of a
point, well inside the reported standard deviation. If the audit stopped here it would be a
reproduction study with a happy ending.

It does not stop here, because that configuration uses a random class-balanced split, and the paper's
footnote defers the split to the code. \tabref{tab:split} changes one flag --- from a random
class-balanced split to the fixed-split path --- and everything else stays identical: same code, same
hyper-parameters, same seed, same epochs. The result is $74.0$ instead of $84.7$.

Ten points is larger than the margin between most of the methods the field compares. A reader who
takes $84.8$ from the paper and sets it beside a public-split baseline is not comparing two methods;
they are comparing two tasks. Nothing in the paper says so, and nothing in the code says so either
--- the flag is one of nineteen in a shell script.

\subsection{SpikeGCL: two reproductions}
\label{sec:gcl}

The third pipeline reproduces on both datasets we ran, and that is what makes the audit
interpretable: the disagreements above are not our inability to run code.

On Cora, with the command in its README, the released code reaches $86.6$ against a published
$87.4\pm0.6$; on Citeseer it reaches $77.3$ against $77.6\pm0.6$. The Citeseer difference
($-0.3$ points) sits inside the published $\pm 0.6$ spread; the Cora difference ($-0.8$ points)
does not. Re-running the pipeline at four seeds gives a measured spread of $0.13$ points on Cora
and $1.24$ points on Citeseer, and a same-nominal-seed re-run of the Citeseer cell differed by
$1.0$ points from the originally audited run --- recorded as part of the sensitivity evidence
rather than reconciled away. The comparison against the published $\pm 0.6$ is the only spread
basis used here; an earlier draft attributed a spread measured on SpikingGCN's Cora seeds
($0.48$--$0.53$) to this pipeline, which was a mis-attribution, and quoted a PubMed range that no
run of this pipeline produced (SpikeGCL's paper reports no PubMed cell).

\subsection{The same baseline, two splits, no label}
\label{sec:mixed}

The last finding is not about a pipeline we ran but about the tables we read while choosing what to
run, and it explains the discrepancy that motivated the audit.

SpikingGCN's paper prints \emph{both} of its splits for every method it compares, and defines them
in its own words: ``In this split method (i.e., Split I), 20 instances from each class are sampled
as the training datasets'' and ``For the second data split (i.e., Split II), the ratio of training
to testing is 8:2, and 20\% of training samples is further used for validation''
\citep{zhu2022spiking}. The 500 validation and 1000 test nodes of Split~I are the standard
Planetoid sizes their loader leaves untouched, and the sampling is visible in their own
\texttt{sample\_per\_class(20, $\cdot$)} call. The two splits are not close. For SpikingGCN itself
the difference is $8.0$ points on Cora, $5.1$ on Citeseer and $11.7$ on PubMed --- larger than the
spread between the top methods in any of the field's comparison tables.

Against that, the two later papers we audited cite the same baseline on different splits, and the
inference is anchored on the quoted definitions rather than on paraphrase. SpikeGraphormer's split
is defined in code (their \texttt{run.sh} sets \texttt{-{-}rand\_split\_class -{-}valid\_num 500
-{-}test\_num 1000} for the published setting and \texttt{-{-}protocol semi} for the public split);
its table lists GCN at $81.5$ on Cora, which is consistent with the Split~I family ($81.0$).
SpikeGCL's table lists SpikingGCN at $87.7$ on Cora, $77.7$ on Citeseer and $87.5$ on PubMed ---
consistent, to within a point, with SpikingGCN's \emph{Split~II} numbers ($88.7$, $77.6$, $89.3$),
not its Split~I numbers ($80.7$, $72.5$, $77.6$). We phrase both readings as ``consistent with'':
the tables do not label the split, and the point is the missing label.

Neither table says which split it is quoting. A reader assembling a baseline table from these two
papers therefore mixes the splits without being told, and the resulting table carries a systematic
error of five to twelve points on some rows and not others. This is not a citation mistake in the
ordinary sense --- both numbers are correct and both papers report what they measured --- and that is
precisely why it has survived: there is no error to find, only an unlabelled choice.

\subsection{The population, not a sample}
\label{sec:census}


A three-method audit would be a sample if more methods were available. They are not.
\tabref{tab:census} lists the seven spiking-GNN methods in our related work and whether each declares
a code repository. Three do; we audited all three. The remaining four publish no code in their arXiv
record or their paper, so they cannot be audited by us or by anyone else, and any number attributed
to them in a comparison table is unverifiable in principle.

This changes what the audit is. It is not "three methods out of many"; within the seven-method
survey that frames this paper --- the spiking-GNN methods collected in our related work from the
2021--2026 literature, with repository status checked against each arXiv record and paper in
November 2026 --- the three auditable methods are the complete auditable set of that survey at the
time of writing. The claim is deliberately bounded: it does not assert a field-wide census, and the
reason the auditable set is small is itself part of the finding.

\subsection{What can now be placed side by side}

With three pipelines measured on the public split, four numbers become commensurable on Cora:
$50.0$ (SpikingGCN, its own code), $74.0$ (SpikeGraphormer, its own code with the split changed),
$86.6$ (SpikeGCL, self-supervised, its own code) --- and, shown for context and not counted as an
audited row, $80.5$ from the conversion study of our companion manuscript \citep{companion}, whose
own registered result files carry the value and whose authors' scripts are not part of this audit.
The published ordering --- SpikingGCN $80.7$, then SpikeGraphormer
$84.8$, then SpikeGCL $87.4$ --- is not the measured ordering, and the largest single move belongs to
the method that reproduces most exactly.

The four numbers measure three training paradigms on one dataset, so a ranking across paradigms
would be a category error. What they establish is narrower: a baseline table assembled from the
literature would have been wrong about all four.
\section{Discussion}
\label{sec:discussion}

\paragraph{Three failures with three different remedies.} The audit's value is that it separates
things that are usually lumped together under "reproducibility".

\emph{A number that reproduces on some datasets and misses on others} is harder to diagnose than one
that never reproduces, because the environment is demonstrably working: the same SpikingGCN script
reaches the published CiteSeer value and misses the Cora value by thirty points. The remedy is a
per-dataset reporting convention and a released configuration, not a global environment fix.

\emph{A number that reproduces on a split nobody else uses} is more insidious, because nothing is
broken. The code is correct, the script is honest, the number is real; it simply measures a different
task from the one the reader will compare it with, and the difference is ten points. The remedy is not
more code but a reporting convention: a table entry should carry its split, its number of runs, and
whether the split is the standard one, in the table rather than in a footnote.

\emph{An unlabelled split choice} is the most damaging, because it corrupts a comparison table
without anyone making a mistake. SpikingGCN's own paper prints $80.7$ on Split~I and $88.7$ on
Split~II; the later tables cite $81.5$ and $87.7$ without saying which. Both numbers are correct.
The remedy is a reporting convention: a baseline row should name its split, and re-running a baseline
should be the default rather than the exception.

\paragraph{What a reader should do.} Practically: do not assemble a baseline table from published
numbers in this area. Pick the method you intend to compare against, run its code on your split, and
report that. Each pipeline takes one working session to re-run once its environment is pinned,
which is what makes the recommendation practical.

\paragraph{Why the field is unusually exposed.} Two structural facts make spiking GNNs worse than the
average subfield. The first is that the spiking encoder is a stochastic sampler reading its input as a
probability, so a preprocessing choice that is cosmetic in a conventional GNN --- row-normalising the
features --- silently disables learning rather than degrading it. The second is that a spiking
network's accuracy depends on the number of timesteps, a hyper-parameter that is reported as part of
the method's efficiency claim rather than as part of its accuracy claim, so two numbers can differ by
ten points without either paper being wrong about what it measured.

\paragraph{Division of evidence with the companion paper.} The conversion study \citep{companion}
uses two numbers from this audit --- $74.0$ and $86.6$ on the public split --- as an external anchor
for its cost comparison. It does not report the audit protocol, the reproduction grid or the
split-mixing analysis, and none of its claims depends on them. This paper reports those results and
does not use the conversion method's fidelity, cost or theory results. The overlap between the two
manuscripts is therefore two measured numbers, explicitly cross-cited, rather than a shared result
set.

\paragraph{Threats to this audit.} Ours, not the audited papers'. \emph{Environment drift}: where a
stated dependency no longer exists for our platform we installed the nearest available version, and
recorded it; a reader with the original stack may get different numbers, and if they do, that is
itself the point. \emph{Under-powered cells}: four of the seven rows are single-run, and we draw no
conclusion from a difference smaller than the spread the multi-seed rows show. \emph{Cora-centred}:
four of seven rows are Cora, chosen because it is the dataset all three papers report, not because it
is representative; a dataset where all three methods agree might behave differently. \emph{Our own
code is in the comparison}: the $80.5$ in \secref{sec:results} comes from a companion paper, and a
reader should discount it accordingly; it is not load-bearing for any claim here.

\paragraph{Evaluative claims and their scope.} The audit supports one claim: a published accuracy in
this area is not transportable to a different split or codebase without being re-run. It does not
support a ranking of the three methods, a judgement on the correctness of any audited paper, or a
claim about spiking GNNs beyond the citation graphs measured here. That claim rests on three
assumptions --- that the released repository is the artefact the paper describes, that the pinned
environment reproduces the authors' execution conditions, and that the numbers in a comparison table
belong to the methods they name. Each is stated because each is checkable, and \secref{sec:mixed} is
what happens when the third fails.

\paragraph{Broader impact and intended use.} An executable audit is easy to misread as a quality
ranking, so every reproduced value here carries its split, its run count and the configuration that
produced it, and no audited paper is described as wrong. The constructive use is prospective: a new
method can be compared against a number produced on its own split by running the released code,
instead of against a figure copied from a table. This study introduces no new dataset --- all three
graphs are the public Planetoid splits used by the audited papers --- and the released artefact is a
protocol with its logs rather than a data resource.

\paragraph{What would falsify us.} Each of the three findings is a claim that a reader can check
against shipped artefacts. If SpikingGCN's code produces $80.7$ on Cora for you, our environment is
wrong and we would like to know. If SpikeGraphormer's public-split configuration produces $84.8$ for
you, we have mis-set a flag. If a later table labels which split it is quoting, the split-mixing
finding no longer applies to that table.
\section{Conclusion}
\label{sec:conclusion}

We ran every spiking graph neural network in our related work that publishes code --- three of the
seven surveyed methods --- using each pipeline's own entry point with its configuration and
environment (disclosing the compatibility patches two of the three required), and reported what came
out.

One reproduces on some datasets and not others: SpikingGCN's code gives $50.0$ on Cora where its
paper reports $80.7$, $64.3$ on PubMed against $77.6$, and $70.6$ on CiteSeer against $72.5$. The
CiteSeer result rules out a broken environment as the explanation for the Cora gap; the input
diagnosis localises what is not the cause (self-loops are not decisive at $51.5$ either way) and
what is (row-normalisation collapses every variant to $13.0$, unrescued by the conventional
rescaling fix). One reproduces exactly and is still unusable as a baseline:
SpikeGraphormer's $84.8$ is a random-split number, and the same code on the public split gives
$74.0$. One reproduces on two datasets --- and the tables we read to choose it cite one baseline on
two different splits, Split~I and Split~II, without labelling them, although those splits differ by
five to twelve points.

These are three different failures, and the fact that they are different is the finding. A field that
worries about variance and seed sensitivity would not catch any of them: every number involved is
internally consistent and was produced by the code that claims to produce it. The transportability
verdict that replaces the variance question is per dataset: measured against seed spreads, the split
gap dominates noise on Cora ($20.5\times$ our pre-registered $3\sigma$ bar; $61.5\times$ the seed
s.d.), is marginal on CiteSeer ($1.37\times$ the bar), and is confounded with seed noise on PubMed
($0.73\times$ the bar), where the two effects cannot be separated at the measured seed count.

The practical recommendation is a single sentence: in this area, do not assemble a baseline table from
published numbers. Run the code. We have made that cheap by shipping the environments, the commands
and the logs, and by writing down which of our own deviations a reader should discount.

The audit covers the complete auditable set of our seven-method survey: three methods. The other four
publish no code that could be checked; the census claim is bounded to that survey and its November
2026 record checks, not to the field.

\bibliographystyle{plainnat}
\bibliography{refs}

\end{document}